% ============================================================
% 第1章补充：经典例题与自学元素
% 插入位置：第1章现有miniquiz之后（\end{miniquiz} 之后）
% ============================================================

\section*{习题精选与详解}
\addcontentsline{toc}{section}{习题精选}

\begin{example}{柱坐标下的散度与旋度}{cylindrical-ops}
计算柱坐标 $(r,\theta,z)$ 下矢量场 $\vect{A} = r\unitvec{r} + z\unitvec{z}$ 的散度与旋度。
\tcblower
\textbf{解答：}
柱坐标散度公式：
\[
\divg\vect{A} = \frac{1}{r}\frac{\partial(rA_r)}{\partial r} + \frac{1}{r}\frac{\partial A_\theta}{\partial\theta} + \frac{\partial A_z}{\partial z}
\]
代入 $A_r = r,\; A_\theta = 0,\; A_z = z$：
\[
\divg\vect{A} = \frac{1}{r}\frac{\partial(r^2)}{\partial r} + 0 + \frac{\partial z}{\partial z} = \frac{1}{r}\cdot 2r + 1 = 2 + 1 = 3
\]

柱坐标旋度公式：
\[
\curl\vect{A} = \left(\frac{1}{r}\frac{\partial A_z}{\partial\theta} - \frac{\partial A_\theta}{\partial z}\right)\unitvec{r} + \left(\frac{\partial A_r}{\partial z} - \frac{\partial A_z}{\partial r}\right)\unitvec{\theta} + \frac{1}{r}\left(\frac{\partial(rA_\theta)}{\partial r} - \frac{\partial A_r}{\partial\theta}\right)\unitvec{z}
\]
逐项计算：
\begin{itemize}[leftmargin=*,itemsep=0pt]
\item $r$ 分量：$\frac{1}{r}\frac{\partial z}{\partial\theta} - \frac{\partial 0}{\partial z} = 0 - 0 = 0$
\item $\theta$ 分量：$\frac{\partial r}{\partial z} - \frac{\partial z}{\partial r} = 0 - 1 = -1$
\item $z$ 分量：$\frac{1}{r}\left[\frac{\partial(0)}{\partial r} - \frac{\partial r}{\partial\theta}\right] = 0 - 0 = 0$
\end{itemize}
故
\[
\curl\vect{A} = -\unitvec{\theta}
\]
\textbf{要点：} 虽然 $A_r$ 和 $A_z$ 均不显含 $\theta$，但 $A_r=r$ 导致 $\theta$ 分量非零。柱坐标下旋度的 $r$、$\theta$ 项容易遗漏 $1/r$ 因子，需特别留意。
\end{example}

\begin{example}{麦克斯韦分布的矩与平均动能}{ch1-maxwell-moments}
设一维麦克斯韦速度分布为
\[
f(v_x) = n\sqrt{\frac{m}{2\pi k_{\!B}T}}\,\exp\!\left(-\frac{mv_x^2}{2k_{\!B}T}\right)
\]
求平均速度 $\langle v_x\rangle$、均方速度 $\langle v_x^2\rangle$ 以及平均动能 $\langle \frac{1}{2}mv_x^2\rangle$。
\tcblower
\textbf{解答：}
由对称性，$f(v_x)$ 是偶函数，被积函数 $v_x f(v_x)$ 是奇函数，因此
\[
\langle v_x\rangle = \frac{1}{n}\int_{-\infty}^{+\infty} v_x f(v_x)\,\diff v_x = 0
\]

计算 $\langle v_x^2\rangle$，利用高斯积分公式
\[
\int_{-\infty}^{+\infty} x^2 \,\mathrm{e}^{-ax^2}\diff x = \frac{\sqrt{\pi}}{2a^{3/2}}
\]
令 $a = \frac{m}{2k_{\!B}T}$，则
\[
\langle v_x^2\rangle = \sqrt{\frac{m}{2\pi k_{\!B}T}}\int_{-\infty}^{+\infty} v_x^2 \,\mathrm{e}^{-av_x^2}\diff v_x = \sqrt{\frac{m}{2\pi k_{\!B}T}}\cdot\frac{\sqrt{\pi}}{2a^{3/2}}
\]
代入 $a = \frac{m}{2k_{\!B}T}$：
\[
\langle v_x^2\rangle = \sqrt{\frac{m}{2\pi k_{\!B}T}}\cdot\frac{\sqrt{\pi}}{2}\left(\frac{2k_{\!B}T}{m}\right)^{3/2} = \frac{k_{\!B}T}{m}
\]

平均动能：
\[
\left\langle\frac{1}{2}mv_x^2\right\rangle = \frac{1}{2}m\cdot\frac{k_{\!B}T}{m} = \frac{1}{2}k_{\!B}T
\]
\textbf{推广：} 三维麦克斯韦分布中，每个自由度贡献 $\frac{1}{2}k_{\!B}T$，总平均动能为 $\frac{3}{2}k_{\!B}T$。
\end{example}

\begin{example}{$\delta$函数与高斯函数的傅里叶变换}{fourier-transforms}
利用本书傅里叶变换定义
\[
\tilde{f}(k) = \int_{-\infty}^{+\infty} f(x)\,\mathrm{e}^{-\mathrm{i} kx}\,\diff x
\]
求：\textbf{(1)} $\delta(x)$ 的傅里叶变换；\textbf{(2)} $f(x)=\exp\!\left(-\frac{x^2}{2\sigma^2}\right)$ 的傅里叶变换。
\tcblower
\textbf{解答：}

\textbf{(1)} 对 $\delta(x)$，利用筛选性质：
\[
\tilde{\delta}(k) = \int_{-\infty}^{+\infty} \delta(x)\,\mathrm{e}^{-\mathrm{i} kx}\,\diff x = \mathrm{e}^{-\mathrm{i} k\cdot 0} = 1
\]
物理意义：空间上无限窄的脉冲包含所有波数分量，且各分量振幅相等。

\textbf{(2)} 对高斯函数：
\[
\tilde{f}(k) = \int_{-\infty}^{+\infty} \exp\!\left(-\frac{x^2}{2\sigma^2}\right)\mathrm{e}^{-\mathrm{i} kx}\,\diff x
\]
对指数配方：
\[
-\frac{x^2}{2\sigma^2} - \mathrm{i}kx = -\frac{1}{2\sigma^2}\left(x^2 + 2\mathrm{i}k\sigma^2 x\right) = -\frac{1}{2\sigma^2}\left[(x+\mathrm{i}k\sigma^2)^2 + k^2\sigma^4\right]
\]
利用高斯积分（复平移不改变积分值）：
\[
\tilde{f}(k) = \exp\!\left(-\frac{k^2\sigma^2}{2}\right)\int_{-\infty}^{+\infty} \exp\!\left(-\frac{u^2}{2\sigma^2}\right)\diff u = \sqrt{2\pi}\,\sigma\,\exp\!\left(-\frac{k^2\sigma^2}{2}\right)
\]
\textbf{结论：} 高斯函数的傅里叶变换仍为高斯函数；空间域越窄（$\sigma$ 越小），波数域越宽（$k$ 展布越宽），体现了量子/波动中的不确定性关系。
\end{example}

\begin{example}{球坐标下的梯度与拉普拉斯}{spherical-gradient}
求标量场 $f(r)=\frac{1}{r}$ 在球坐标 $(r,\theta,\varphi)$ 下的梯度，并验证其满足拉普拉斯方程 $\lap f=0$（$r\neq 0$）。
\tcblower
\textbf{解答：}
球坐标梯度公式：
\[
\grad f = \frac{\partial f}{\partial r}\unitvec{r} + \frac{1}{r}\frac{\partial f}{\partial\theta}\unitvec{\theta} + \frac{1}{r\sin\theta}\frac{\partial f}{\partial\varphi}\unitvec{\varphi}
\]
由于 $f$ 只依赖于 $r$，后两项为零：
\[
\grad f = \frac{\partial}{\partial r}\left(\frac{1}{r}\right)\unitvec{r} = -\frac{1}{r^2}\unitvec{r}
\]

验证拉普拉斯（$r\neq 0$）。球坐标拉普拉斯：
\[
\lap f = \frac{1}{r^2}\frac{\partial}{\partial r}\!\left(r^2\frac{\partial f}{\partial r}\right) + \frac{1}{r^2\sin\theta}\frac{\partial}{\partial\theta}\!\left(\sin\theta\frac{\partial f}{\partial\theta}\right) + \frac{1}{r^2\sin^2\theta}\frac{\partial^2 f}{\partial\varphi^2}
\]
后两项为零，第一项：
\[
\lap f = \frac{1}{r^2}\frac{\partial}{\partial r}\!\left[r^2\cdot\left(-\frac{1}{r^2}\right)\right] = \frac{1}{r^2}\frac{\partial}{\partial r}(-1) = 0
\]
\textbf{说明：} $f=1/r$ 是三维拉普拉斯方程的基本解（除原点外），在静电学中对应点电荷的势。原点处的奇异性由 $\delta$ 函数描述：$\lap(1/r)=-4\pi\delta(\vect{r})$。
\end{example}

% ch1 新例题：分布函数矩的计算
\begin{example}{麦克斯韦分布的矩与最概然速率}{ch1-maxwell-dist}
一维麦克斯韦分布 $f(v) = \dfrac{n}{\sqrt{\pi}\,v_t}\exp(-v^2/v_t^2)$，其中 $v_t = \sqrt{2k_{\!B}T/m}$（\textbf{注意}：此处 $v_t$ 为一维分布的尺度参数，满足 $v_t = \sqrt{2}\,v_{\mathrm{rms}}$，其中 $v_{\mathrm{rms}}=\sqrt{k_{\!B}T/m}$ 为一维方均根速率；不要与三维最概然速率 $v_p=\sqrt{2k_{\!B}T/m}$ 的数值巧合混淆——两者表达式相同但含义不同）。计算：
（1）平均速率 $\langle v\rangle$；（2）方均根速率 $v_{\mathrm{rms}}=\sqrt{\langle v^2\rangle}$；（3）最概然速率 $v_p$，并比较三者的相对大小。
\tcblower
\textbf{解答：}

\textbf{(1) 平均速率：}
\[
\langle v\rangle = \int_{-\infty}^{+\infty} v\,f(v)\,\diff v = 0
\]
被积函数 $v\,f(v)$ 是奇函数，对称区间积分为零——一维速度的平均值为零（无净流动）。

\textbf{(2) 方均根速率：}
\[
\langle v^2\rangle = \int_{-\infty}^{+\infty} v^2\,f(v)\,\diff v = \frac{n}{\sqrt{\pi}v_t}\cdot v_t^3\int_{-\infty}^{+\infty} x^2\ee^{-x^2}\,\diff x = \frac{n v_t^2}{\sqrt{\pi}}\cdot\frac{\sqrt{\pi}}{2} = \frac{n v_t^2}{2}
\]
按密度归一化后 $\langle v^2\rangle = v_t^2/2 = k_{\!B}T/m$（一维能量均分定理，每自由度 $k_{\!B}T/2$）：
\[
v_{\mathrm{rms}} = \sqrt{\frac{k_{\!B}T}{m}} = \frac{v_t}{\sqrt{2}}
\]

\textbf{(3) 最概然速率：} 令 $f'(v)=0$：
\[
\frac{\mathrm{d}f}{\mathrm{d}v} = -\frac{2v}{v_t^2}f(v) = 0 \quad\Rightarrow\quad v_p = 0
\]
一维情形最概然速率为零（分布峰值在 $v=0$）。

\textbf{对比与讨论：} 三维情形下三者不相等：$v_p : \bar{v} : v_{\mathrm{rms}} = 1 : 1.128 : 1.225$。分布不对称（右偏），故 $\langle v\rangle > v_p$；均方速率对高速尾部敏感，故 $v_{\mathrm{rms}} > \langle v\rangle$。在等离子体物理中 $v_{\mathrm{rms}}$（即本书的热速度定义之一）最常用，因为它直接与温度通过能量均分定理联系。
\end{example}

\begin{formulabox}{第1章核心公式速查}{ch1-formulas}
\begin{itemize}[leftmargin=*,itemsep=0pt]
\item 柱坐标散度：$\displaystyle\divg\vect{A} = \frac{1}{r}\frac{\partial(rA_r)}{\partial r} + \frac{1}{r}\frac{\partial A_\theta}{\partial\theta} + \frac{\partial A_z}{\partial z}$
\item 柱坐标旋度：$\displaystyle\curl\vect{A} = \left(\frac{1}{r}\frac{\partial A_z}{\partial\theta} - \frac{\partial A_\theta}{\partial z}\right)\unitvec{r} + \left(\frac{\partial A_r}{\partial z} - \frac{\partial A_z}{\partial r}\right)\unitvec{\theta} + \frac{1}{r}\left(\frac{\partial(rA_\theta)}{\partial r} - \frac{\partial A_r}{\partial\theta}\right)\unitvec{z}$
\item 球坐标梯度：$\displaystyle\grad f = \frac{\partial f}{\partial r}\unitvec{r} + \frac{1}{r}\frac{\partial f}{\partial\theta}\unitvec{\theta} + \frac{1}{r\sin\theta}\frac{\partial f}{\partial\varphi}\unitvec{\varphi}$
\item 球坐标拉普拉斯：$\displaystyle\lap f = \frac{1}{r^2}\frac{\partial}{\partial r}\!\left(r^2\frac{\partial f}{\partial r}\right) + \frac{1}{r^2\sin\theta}\frac{\partial}{\partial\theta}\!\left(\sin\theta\frac{\partial f}{\partial\theta}\right) + \frac{1}{r^2\sin^2\theta}\frac{\partial^2 f}{\partial\varphi^2}$
\item 分布函数矩：$\displaystyle\langle g(\vect{v})\rangle = \frac{1}{n}\int g(\vect{v})f(\vect{v})\,\diff^3v$
\item 平均动能（麦克斯韦）：$\displaystyle\frac{1}{2}m\langle v^2\rangle = \frac{3}{2}k_{\!B}T$
\item 傅里叶变换：$\displaystyle\tilde{f}(k) = \int_{-\infty}^{+\infty} f(x)\,\mathrm{e}^{-\mathrm{i} kx}\,\diff x$
\item 逆变换：$\displaystyle f(x) = \frac{1}{2\pi}\int_{-\infty}^{+\infty} \tilde{f}(k)\,\mathrm{e}^{+\mathrm{i} kx}\,\diff k$
\item $\delta$ 函数筛选性：$\displaystyle\int_{-\infty}^{+\infty} f(x)\delta(x-a)\,\diff x = f(a)$
\end{itemize}
\end{formulabox}

\begin{checklistbox}{第1章自学检查清单}{ch1-checklist}
\begin{itemize}[leftmargin=*,label=$\square$]
\item 我能在柱坐标和球坐标下正确写出 $\grad$、$\divg$、$\curl$ 和 $\lap$
\item 我理解了分布函数 $f(\vect{v})$ 的物理意义，并能计算密度、平均速度、压强
\item 我掌握了傅里叶变换和逆变换的公式，并知道本书的符号约定（指数项 $-\mathrm{i}kx$）
\item 我能用 $\delta$ 函数表示点电荷和点粒子的密度分布
\item 我了解二阶张量的基本运算（并矢、缩并、迹）
\end{itemize}
\end{checklistbox}

\endinput
